\subsection{Optimality of the t-exact realisation}

The construction above has a converse when the marked kernel embedding is required to preserve the standard t-structures.  The converse applies to arbitrary ordinary homological ring epimorphisms, without a projectivity assumption on the target.

\begin{theorem}[T-exact kernel realisations are two-term]\label{thm:texact-optimality}
Let \(\lambda:R\to S\) be a homological ring epimorphism of ordinary rings and suppose
\[
 \cK_\lambda=\ker(S\Ltensor_R-)=\Loc(T),
 \qquad T\in K^b(R\text{-}\mathrm{proj}).
\]
If \(B=\RHom_R(T,T)^{\op}\) is connective and its Morita embedding
\(F:D(B)\to D(R)\) is t-exact for the standard t-structures, then \(T\) is homotopy equivalent to a complex of finitely generated projectives concentrated in degrees \([-1,0]\).  In particular \(H^i(B)=0\) for \(i<-1\).
\end{theorem}

\begin{proof}
A fully faithful t-exact functor reflects standard cohomological bounds.  Indeed, applying it to a truncation triangle gives a truncation triangle, whence
\(H^n(FY)\cong F_{\heartsuit}(H^nY)\); the fully faithful heart functor reflects zero objects.  This argument does not require \(Y\) to be bounded.

Homologicality gives the localisation triangle
\[
 K_R\longrightarrow R\xrightarrow{\lambda}S\longrightarrow K_R[1],
 \qquad K_R=\Cone(R\to S)[-1]\in\cK_\lambda.
\]
To check the last assertion in all degrees, apply \(S\Ltensor_R-\): the resulting map \(S\to S\Ltensor_RS\) is the inverse of the multiplication isomorphism.  With our cohomological convention, the displayed representative of \(K_R\) has terms in degrees \(0,1\).  Thus, writing \(K_R\cong FC\), reflection gives \(C\in D^{[0,1]}(B)\).  For every \(n\), adjunction gives
\[
 \Hom_{D(R)}(T,S[n])
 \cong\Hom_{D(S)}(S\Ltensor_RT,S[n])=0.
\]
Applying \(\RHom_R(T,-)\) to the triangle identifies the cohomology of \(T^\vee=\RHom_R(T,R)\) with that of \(\RHom_R(T,K_R)\).  Full faithfulness and \(F(B)=T\) identify the latter groups with \(\Hom_{D(B)}(B,C[n])=H^n(C)\).  Hence
\[
 T^\vee\in D^{[0,1]}(R^{\op}),
 \qquad T=F(B)\in D^{\le0}(R).
\]
Only cohomological bounds are asserted here; the two complexes have different additional module actions.

For a bounded projective complex whose cohomology vanishes above degree \(b\), the differential onto its highest term above \(b\) is a split surjection.  Removing the associated contractible summand and repeating eliminates all terms above \(b\).  Apply this first to \(T\), obtaining a projective representative with no positive terms.  Its finite-projective dual has no negative terms.  Apply the same cancellation to that dual above degree one; it is then supported in \([0,1]\).  Finite-projective biduality gives a representative of \(T\) in \([-1,0]\).  All steps preserve homotopy equivalence.  Its endomorphism complex has terms only in degrees \([-1,1]\), proving the vanishing below \(-1\).
\end{proof}

\begin{corollary}[Universality of a t-exact realisation]\label{cor:texact-universal}
Under \Cref{thm:texact-optimality}, write a two-term representative as \(T=[P^{-1}\xrightarrow{\sigma}P^0]\).  Then \(\lambda\) is the universal localisation at \(\sigma\).
\end{corollary}

\begin{proof}
The vanishing \(S\Ltensor_RT=0\) says that \(S\otimes_R\sigma\) is invertible.  The essential image of derived restriction is \(\cK_\lambda^\perp\).  A stalk \(R\)-module \(M\) belongs to this image exactly when \(\Hom_R(\sigma,M)\) is invertible.  Its preimage is a stalk \(S\)-module, because restriction is exact and reflects cohomological vanishing.  Its extending \(S\)-action is unique: full faithfulness applied to the identity of the underlying \(R\)-module identifies any two such actions.

Let \(f:R\to C\) invert \(\sigma\).  Tensor--Hom adjunction puts \({}_RC\) in \(\cK_\lambda^\perp\), giving its unique extending \(S\)-action.  Every right multiplication on \(C\) is \(R\)-linear and therefore \(S\)-linear.  Consequently \(s\cdot c=(s\cdot1)c\), and \(g(s)=s\cdot1\) is a unital ring homomorphism \(S\to C\) extending \(f\).  Any other extension defines the same \(S\)-action, hence equals \(g\).  This is the required universal property.
\end{proof}

\begin{corollary}[Classification of t-exact marked realisations]\label{cor:texact-classification}
For a proper connective dg algebra \(B\) over a field, the following are equivalent:
\begin{enumerate}[label=\textup{(\roman*)}]
\item \(B\) is quasi-isomorphic to the dg endomorphism algebra, with opposite, of a perfect generator of the kernel of a homological epimorphism of finite-dimensional algebras, with standard t-exact marked kernel embedding;
\item \(H^i(B)=0\) for \(i<-1\);
\item such a realisation exists as a universal localisation at one map, with target projective on both sides and full bounded profile.
\end{enumerate}
\end{corollary}

\begin{proof}
The implication (i)\(\Rightarrow\)(ii) is \Cref{thm:texact-optimality}; (ii)\(\Rightarrow\)(iii) is \Cref{thm:finite-crossed,thm:realisation}; and (iii)\(\Rightarrow\)(i) is immediate.  Quasi-isomorphisms of connective dg models preserve the standard t-structures.
\end{proof}

Thus allowing ordinary homological epimorphisms which are not assumed to be universal localisations does not enlarge this t-exact realisation range.  The full-profile seed of \Cref{thm:finite-crossed} still realises every cut by \Cref{thm:ambient}; an extension realising a proper cut need not retain the seed's projectivity or t-exactness.

\subsection{Projective targets and the Schur boundary}

\begin{corollary}[Kernel categories with projective right target]\label{cor:projective-target-classes}
The dg Morita classes of compact kernel categories of homological epimorphisms of finite-dimensional algebras with projective right target are exactly the classes admitting a proper connective dg model with cohomology in \([-1,0]\).  Every such class has a realisation with target projective on both sides, standard t-exact marking by a thin model, and full bounded profile.
\end{corollary}

\begin{proof}
Let \(S_R\) be projective.  By \Cref{prop:left-projection}, the kernel has the perfect generator
\[
 G=\Cone\!\left(\Hom_{R^{\op}}(S,R)\longrightarrow R\right),
\]
with terms in \([-1,0]\) and \(\Hom(G,G[n])\cong H^n(G)\).  Thus \(\RHom_R(G,G)^{\op}\) is proper with cohomology in \([-1,0]\), and models the compact kernel.  Conversely apply \Cref{thm:finite-crossed,thm:realisation} to a thin model.  A dg Morita equivalence preserves \(\Dfd\), by its intrinsic finite-total-Hom characterisation, so the full-profile assertion is independent of the compact-generator marking.  The t-exact assertion concerns the thin marking supplied by the construction.
\end{proof}

\begin{proposition}[Negative Calabi--Yau obstruction]\label{prop:negative-cy-obstruction}
Let \(B\) be a proper dg algebra such that \(\per B\ne0\) has Serre functor \([-d]\), where \(d\ge2\).  Then \(\per B\) is not equivalent to the compact kernel of either
\begin{enumerate}[label=\textup{(\roman*)}]
\item a homological epimorphism of finite-dimensional algebras with \(S_R\in\per(R^{\op})\); or
\item a finite stably flat universal localisation of finite-dimensional algebras, without any perfection assumption on the target.
\end{enumerate}
Neither exclusion requires t-exactness or full profile.
\end{proposition}

\begin{proof}
Every nonzero \(G\in\per B\) satisfies
\[
 \Hom(G,G[-d])\cong D_k\End(G)\ne0
\]
by Serre duality.  In (i), \Cref{prop:left-projection} supplies a nonzero perfect generator with \(\Hom(G,G[n])=0\) for \(n<-1\), a contradiction.  In (ii), the transported sum of the presentation cones has the same vanishing by \Cref{prop:thin}, again a contradiction.
\end{proof}

\begin{example}[Two excluded formal algebras]\label{ex:negative-cy-algebras}
The dg algebras
\[
 B_2=k[u]/(u^2),\quad |u|=-2,
 \qquad
 B_{\wedge}=\Lambda_k(u,v),\quad |u|=|v|=-1,
\]
with zero differential satisfy \Cref{prop:negative-cy-obstruction} with \(d=2\), over every field.  For \(B_2\), take the trace with value one on \(u\) and zero on \(1\); for \(B_{\wedge}\), take the coefficient of \(uv\).  The Gram matrices of \((a,b)\mapsto\operatorname{tr}(ab)\), in the respective bases \((1,u)\) and \((1,u,v,uv)\), are
\[
 \begin{pmatrix}0&1\\1&0\end{pmatrix},
 \qquad
 \begin{pmatrix}0&0&0&1\\0&0&1&0\\0&-1&0&0\\1&0&0&0\end{pmatrix}.
\]
These pairings are nondegenerate and graded symmetric, including in characteristic two, and give dg bimodule isomorphisms \(D_kB\cong B[-2]\).  Finite-projective tensor--Hom duality gives
\[
 D_k\RHom_B(X,Y)\simeq
 \RHom_B\bigl(Y,D_kB\Ltensor_BX\bigr)
 \qquad(X,Y\in\per B).
\]
The formula holds on finite free modules and extends by shifts, cones and direct summands; it identifies \([-2]\) as the Serre functor.
\end{example}

\begin{example}[A Schur realisation outside the epimorphism range]\label{ex:sphere-schur}
Let \(A\) be the radical-square-zero algebra of the oriented three-cycle
\[
 1\xrightarrow{a_1}2\xrightarrow{a_2}3\xrightarrow{a_3}1.
\]
Its basis is \(e_1,e_2,e_3,a_1,a_2,a_3\), with orthogonal vertex idempotents, the source and target identities of the arrows, and all products of two arrows zero.  Put \(e=e_2+e_3\), so \(eAe=k(2\to3)\), and take
\[
 T=[Ae_2\xrightarrow{a_2}Ae_3\xrightarrow{a_3}Ae_1]
 \quad\text{in degrees }[-2,-1,0].
\]
Then there is a recollement
\[
 D(B_2)\longrightarrow D(A)\xrightarrow{e(-)}D(eAe)
\]
whose kernel embedding is t-exact and has full bounded profile.  The source has dimension six and the quotient corner has dimension three.  The quotient functor here is the Schur functor, not extension along an ordinary ring epimorphism.

Indeed, the endomorphism complex of \(T\) has dimensions \(1,0,3,2\) in degrees \(-2,-1,0,1\), and zero in every other degree.  Its only nonzero differential sends the three diagonal scalars to their two consecutive differences and has rank two over every field.  Thus its cohomology is \(k\) in degrees \(-2,0\).  Sending \(u\) to right multiplication by \(a_1\), from \(Ae_1\) to \(Ae_2\), and \(1\) to the identity gives a strict quasi-isomorphism
\[
 B_2\longrightarrow\End_A^\bullet(T)^{\op}.
\]
As right dg \(B_2\)-modules, \(e_1T=B_2\); the other two vertex components are identity-contractible two-term complexes, on which \(u\) acts trivially.  Consequently \(T\Ltensor_{B_2}-\) preserves cohomological degrees and finite-dimensional total cohomology, and reflects the latter.

The object \(T\) belongs to \(\ker e(-)\).  If \(X\) belongs to this kernel, the finite filtration of \(\Hom_A^\bullet(T,X)\) by the three terms of \(T\) has two acyclic associated factors, \(e_2X\) and \(e_3X\), leaving \(e_1X\).  Hence \(\RHom_A(T,X)=0\) implies \(X=0\).  The \(\Loc(T)\)-approximation triangle inside \(D(A)\) now gives \(\ker e(-)=\Loc(T)\): its residual object belongs to the kernel and is right-orthogonal to \(T\), so is zero.  Dg Morita theory and the idempotent recollement of \cite[Proposition~6.1]{KalckYangII} yield the displayed recollement.  The complete finite computation is included in the supplement.  By \Cref{thm:texact-optimality,prop:negative-cy-obstruction}, this example cannot be replaced by any of the ordinary-epimorphism realisations excluded there.
\end{example}

\begin{example}[Full profile does not force a thin marking]\label{ex:marking-boundary}
Let \(A=k(1\to2)\), \(R=A\times k\), and \(\lambda:R\to k\) be projection.  The target is projective on both sides and multiplication \(k\otimes_Rk\to k\) is an isomorphism, so all positive Tor groups vanish and the kernel is \(D(A)\).  Put \(P_i=Ae_i\) and \(T=P_1\oplus P_2[2]\).  The complete Hom calculation is
\[
 \End_A(P_i)=k,\qquad
 \Hom_A(P_1,P_2)=k,\qquad\Hom_A(P_2,P_1)=0;
\]
all higher or lower Hom groups of these stalk projectives vanish.  Therefore \(B=\End_R^\bullet(T)^{\op}\) has zero differential,
\[
 H^0(B)=k\times k,\qquad H^{-2}(B)=k,
\]
and no other cohomology.  Since \(T\) and \(A\) classically generate the same perfect category, the finite-total-Hom characterisation gives full profile for this marking.  It is not t-exact by \Cref{thm:texact-optimality}.  Thus \Cref{cor:projective-target-classes} concerns existence of a thin model in a dg Morita class; it does not assert that every connective marked model is thin.  In particular, the broader \Cref{q:realisation} is not settled by the t-exact classification.
\end{example}
